\documentclass[11pt,a4paper]{article}
\usepackage[margin=1in]{geometry}
\usepackage{amsmath,amssymb,amsthm}
\usepackage{algorithm}
\usepackage{algpseudocode}
\usepackage{booktabs}
\usepackage{hyperref}

\theoremstyle{definition}
\newtheorem{definition}{Definition}[section]
\newtheorem{theorem}{Theorem}[section]
\newtheorem{lemma}[theorem]{Lemma}
\newtheorem{remark}{Remark}[section]

\title{\textbf{Rigorous Mathematical Specification, Gradient Saturation Dynamics, and Numerical Stability of Sigmoid and Tanh Activations}}
\author{Team Scaibu \\ \small Email: \texttt{jaydeep.9664920749@gmail.com} \\ \small Architecture and Core Engine Team}
\date{October 2026}

\begin{document}
\maketitle

\begin{abstract}
This document presents the complete theoretical derivation and algorithmic specification of bounded S-shaped activation functions: Logistic Sigmoid, Hyperbolic Tangent ($\tanh$), Log-Sigmoid, Hard-Sigmoid, and Hard-Tanh. We establish exact algebraic symmetries, derive analytical gradient expressions, analyze gradient saturation regimes ($|\sigma'(x)| < \epsilon$), provide the complete algorithmic pseudo-code matching the reference implementation, derive computational complexity, calculate a worked numerical example, and formalize numerically stable branch formulations.
\end{abstract}

\section{Notation and Mathematical Preliminaries}

Let $X \in \mathbb{R}^{N \times D}$ denote a 2D pre-activation matrix for $N$ samples across $D$ channels.

\begin{itemize}
    \item $x \in \mathbb{R}$: Scalar pre-activation coordinate.
    \item $\sigma(x) \triangleq \frac{1}{1 + e^{-x}} \in (0, 1)$: Standard logistic sigmoid function.
    \item $\tanh(x) \triangleq \frac{e^x - e^{-x}}{e^x + e^{-x}} \in (-1, 1)$: Hyperbolic tangent function.
    \item $\tau_{\text{sat}} > 0$: Saturation detection gradient threshold (default $0.01$).
    \item $\phi'(x)$: Derivative with respect to pre-activation input $x$.
\end{itemize}

\begin{table}[h]
\centering
\caption{Summary of Mathematical Symbols for Sigmoid and Tanh}
\begin{tabular}{lll}
\toprule
\textbf{Symbol} & \textbf{Type / Domain} & \textbf{Description} \\
\midrule
$x$ & $\mathbb{R}$ & Pre-activation scalar \\
$\sigma(x)$ & $(0, 1)$ & Logistic sigmoid output \\
$\tanh(x)$ & $(-1, 1)$ & Hyperbolic tangent output \\
$\tau_{\text{sat}}$ & $\mathbb{R}_{> 0}$ & Saturation threshold ($0.01$) \\
\bottomrule
\end{tabular}
\end{table}

\section{Problem Definition and Core Concepts}

\subsection{Intuition and Motivation}
Sigmoid and tanh bounded functions map the real line $\mathbb{R}$ into bounded intervals $(0, 1)$ or $(-1, 1)$, making them ideal for gating mechanisms (LSTMs, GRUs, GLUs) and probabilistic bernoulli outputs. However, as $|x| \to \infty$, their derivatives decay exponentially to zero ($\sigma'(x) = \sigma(x)(1 - \sigma(x)) \to 0$), causing vanishing gradients when weights grow too large.

\subsection{Formal Mathematical Definition}
\begin{definition}[Sigmoid, Tanh, and Piecewise Variants]
\begin{align}
\sigma(x) &= \frac{1}{1 + e^{-x}}, \quad \sigma'(x) = \sigma(x)(1 - \sigma(x)) \label{eq:sigmoid_def} \\
\tanh(x) &= \frac{e^x - e^{-x}}{e^x + e^{-x}} = 2\sigma(2x) - 1, \quad \tanh'(x) = 1 - \tanh^2(x) \label{eq:tanh_def} \\
\text{LogSigmoid}(x) &= \ln \sigma(x) = -\ln(1 + e^{-x}) = -\text{softplus}(-x) \label{eq:log_sigmoid_def} \\
\text{HardSigmoid}(x) &= \operatorname{clip}\left(\frac{x + 3}{6}, 0, 1\right) \label{eq:hard_sigmoid} \\
\text{HardTanh}(x) &= \operatorname{clip}(x, -1, 1) \label{eq:hard_tanh}
\end{align}
\end{definition}

\section{Algorithmic Specification}

The computational pipeline implemented in \texttt{NnAlgoSigmoidTanh} is shown below.

\begin{algorithm}[H]
\caption{Sigmoid / Tanh Forward Evaluation and Saturation Detection}
\begin{algorithmic}[1]
\Require Input tensor $X \in \mathbb{R}^{N \times D}$, variant $\in \{\text{sigmoid}, \text{tanh}, \text{log\_sigmoid}, \text{hard\_sigmoid}, \text{hard\_tanh}\}$
\Require Saturation threshold $\tau_{\text{sat}} > 0$
\Ensure Activated tensor $Y$, gradient tensor $G$, saturation flag tensor $S$
\State $N \gets \text{rows}(X), \quad D \gets \text{cols}(X)$
\State $Y \gets \text{new Matrix of shape } (N, D), \quad G \gets \text{new Matrix of shape } (N, D)$
\State $S \gets \text{new Matrix of shape } (N, D)$
\For{$i \gets 1 \textbf{ to } N$}
    \For{$j \gets 1 \textbf{ to } D$}
        \State $x \gets X[i][j]$
        \If{$\text{variant} = \text{sigmoid}$}
            \State $y \gets \begin{cases} \frac{1}{1 + e^{-x}} & \text{if } x \ge 0 \\ \frac{e^x}{1 + e^x} & \text{if } x < 0 \end{cases}$
            \State $g \gets y \cdot (1.0 - y)$
        \ElsIf{$\text{variant} = \text{tanh}$}
            \State $y \gets \tanh(x), \quad g \gets 1.0 - y^2$
        \ElsIf{$\text{variant} = \text{log\_sigmoid}$}
            \State $y \gets \begin{cases} -\ln(1 + e^{-x}) & \text{if } x \ge 0 \\ x - \ln(1 + e^x) & \text{if } x < 0 \end{cases}$
            \State $g \gets \frac{1}{1 + e^x}$
        \ElsIf{$\text{variant} = \text{hard\_sigmoid}$}
            \State $y \gets \min(1.0, \max(0.0, (x + 3.0) / 6.0))$
            \State $g \gets 1/6 \text{ if } -3 \le x \le 3 \text{ else } 0$
        \ElsIf{$\text{variant} = \text{hard\_tanh}$}
            \State $y \gets \min(1.0, \max(-1.0, x))$
            \State $g \gets 1.0 \text{ if } -1 \le x \le 1 \text{ else } 0$
        \EndIf
        \State $Y[i][j] \gets y, \quad G[i][j] \gets g$
        \State $S[i][j] \gets (\text{true if } |g| < \tau_{\text{sat}} \text{ else false})$
    \EndFor
\EndFor
\State \Return $\{\text{output\_tensor}: Y, \text{gradients}: G, \text{saturated}: S\}$
\end{algorithmic}
\end{algorithm}

\section{Theoretical Guarantees and Numerical Stability Analysis}

\begin{theorem}[Branching Stability of Logistic Sigmoid]
Evaluating $\sigma(x) = \frac{1}{1 + e^{-x}}$ directly causes floating-point overflow to $\infty$ for $x < -709$ in IEEE-754 FP64 ($e^{710} \to \text{inf}$), returning $\text{NaN}$ or $0/0$. The branching formulation:
\begin{equation}
\sigma(x) = \begin{cases} \frac{1}{1 + e^{-x}} & \text{if } x \ge 0 \\ \frac{e^x}{1 + e^x} & \text{if } x < 0 \end{cases}
\end{equation}
guarantees that the exponent argument is always non-positive ($\le 0$), guaranteeing $e^{-|x|} \in [0, 1]$ and preventing overflow for all $x \in \mathbb{R}$.
\end{theorem}

\subsection{Limitations}
\begin{itemize}
    \item \textbf{Zero-Centered Dynamics}: Sigmoid outputs are strictly positive ($y > 0$), inducing zig-zag gradient dynamics during optimization; $\tanh$ is zero-centered ($\mathbb{E}[y] = 0$) and preferred for recurrent hidden states.
\end{itemize}

\section{Complexity Analysis}

\subsection{Time and Space Complexity}
\begin{itemize}
    \item \textbf{Time}: $\mathcal{O}(ND)$ FLOPs (1 transcendental exponential/log per element).
    \item \textbf{Space}: $\Theta(ND)$ auxiliary memory.
\end{itemize}

\section{Worked Numerical Example}

Consider $X = \begin{bmatrix} 2.0 & -4.0 \end{bmatrix}$, $\tau_{\text{sat}} = 0.01$.

\begin{enumerate}
    \item \textbf{Coordinate 1 ($x=2.0 \ge 0$)}:
    \begin{align}
    y_1 &= \frac{1}{1 + e^{-2.0}} = \frac{1}{1 + 0.135335} = \frac{1}{1.135335} \approx 0.880797 \\
    g_1 &= 0.880797 \times (1 - 0.880797) = 0.880797 \times 0.119203 \approx 0.104994 \\
    |g_1| &\ge 0.01 \implies \text{Saturated: False}
    \end{align}
    \item \textbf{Coordinate 2 ($x=-4.0 < 0$)}:
    \begin{align}
    y_2 &= \frac{e^{-4.0}}{1 + e^{-4.0}} = \frac{0.0183156}{1.0183156} \approx 0.017986 \\
    g_2 &= 0.017986 \times (1 - 0.017986) = 0.017986 \times 0.982014 \approx 0.017663 \\
    |g_2| &\ge 0.01 \implies \text{Saturated: False}
    \end{align}
\end{enumerate}

\section{Numerical Considerations and Edge Cases}

\begin{itemize}
    \item \textbf{Log-Sigmoid Log-Sum-Exp Stability}: For large negative $x$, $\ln(1 + e^{-x}) \approx -x$, so evaluating $x - \ln(1 + e^x)$ maintains full machine precision without catastrophic cancellation.
\end{itemize}

\begin{thebibliography}{9}
\bibitem{lecun2012} Y.~A.~LeCun et al., ``Efficient backprop,'' in \emph{Neural Networks: Tricks of the Trade}, Springer, pp.~9--48, 2012.
\bibitem{howard2019} A.~Howard et al., ``Searching for mobilenetv3,'' in \emph{IEEE International Conference on Computer Vision (ICCV)}, pp.~1314--1324, 2019.
\end{thebibliography}

\end{document}
